\documentclass[11pt]{article}
\usepackage{amsmath,amsfonts,amssymb}
\usepackage{graphicx}
\usepackage{booktabs}
\usepackage{hyperref}
\usepackage{authblk}

% =====================================================================
% REVISED SKELETON — honest ablation / methodology framing.
% Design principle: the CONTRIBUTION is the evaluation protocol and the
% clean negative result, NOT the architecture. Structure follows the
% finding. Every section earns its space by advancing the central claim:
% "no quantum component is load-bearing on real held-out images, and
%  here is a reusable protocol that establishes this rigorously."
% =====================================================================

\title{Do the Quantum Parts Help?\\
A Component-Wise Ablation of Physics-Informed Hybrid\\
Quantum--Classical Image Autoencoders}

\author[1]{First Author}
\author[1]{Second Author}
\affil[1]{Institution, Department}
\date{}

\begin{document}
\maketitle

\begin{abstract}
% Lead with the finding, not the architecture. State the protocol as the
% contribution. Name the leakage post-mortem as a transferable lesson.
We study whether the quantum components of a physics-informed hybrid
quantum--classical image autoencoder (Hamiltonian Vision Kernel, HVK ---
MPS features $\to$ variational quantum circuit $\to$ learnable Heisenberg
energy $\to$ classical decoder) contribute measurably to performance. Using a
resource-matched, multi-seed ablation protocol on held-out CIFAR-10, we find
that freezing the circuit, removing entanglement, dropping the MPS encoder,
disabling the Hamiltonian term, and replacing the circuit with a trivial
classical map all match or exceed the full model. A plain local-linear
classical control ties or beats the entangling model on real images. We further
present a worked post-mortem of a circular ``nonlocal advantage'' benchmark that
manufactured a spurious $R^2\!\approx\!1$ through target--feature leakage, and
distill the diagnostic that exposed it. Our contribution is the evaluation
protocol and a clean, reproducible negative result --- a template for
evidence-grade claims in hybrid quantum machine learning.
\end{abstract}

\section{Introduction}
% Motivation is the CREDIBILITY GAP in hybrid QML, not vision.
Motivate the problem as evaluation, not architecture: hybrid quantum vision
models routinely report advantage without isolating which component is
responsible or matching classical resources. State the gap: component-wise,
resource-matched, leakage-audited evaluation is rare. Preview the three
contributions: (i) an ablation/control protocol; (ii) a clean negative result on
held-out CIFAR; (iii) a reusable leakage diagnostic. State up front that this is
a negative result and why that is valuable.

\paragraph{Contributions.}
\begin{itemize}
  \item A resource-matched, multi-seed, freeze-based ablation protocol that
        isolates the marginal value of each quantum component.
  \item Evidence that, under this protocol, no quantum component of HVK is
        load-bearing on held-out CIFAR-10; matched classical controls tie or win.
  \item A leakage-detection diagnostic, demonstrated on a benchmark of our own
        that it invalidates --- a transferable cautionary methodology.
\end{itemize}

\section{Hamiltonian Vision Kernel: Architecture (Compact)}
% ONE section, not six. Enough to make ablations legible; no more.
% Full circuit definitions go to the appendix.
Give the minimal description needed to read the ablations: patch decomposition +
positional encoding; MPS features; a 6-qubit VQC producing Pauli observables; a
learnable Heisenberg energy regularizer; a classical MLP decoder that is the main
trainable component. Present the variants (HVK1D chain, HVK2D grid, U(1)-symmetric,
IBM probe) in a single comparison table. Defer exact gate sequences to
Appendix~\ref{app:circuits}. Explicitly flag that the decoder carries the trainable
capacity --- this foreshadows the result.

\section{Evaluation Protocol}
% THIS is the paper's methodological core — give it prominence.
\subsection{Ablation and Control Design}
Define the freeze-quantum / freeze-classical isolation, no-entanglement,
no-MPS, no-energy, and classical-replacement ablations, and --- critically --- the
\emph{resource matching} (equal feature width, equal readout parameter count) that
makes quantum-vs-classical comparison fair.
\subsection{Statistical Protocol}
Multi-seed (5+) runs with mean~$\pm$~std; matched-convergence training budget;
report that sub-dB gaps within seed noise are not claims.
\subsection{Leakage Auditing}
State the discipline: audit that regression targets are not linear functions of
features handed to only one model. Present it as a general check, motivated by
\S\ref{sec:leakage}.
\subsection{Datasets and Metrics}
Single-image, multi-image, and held-out CIFAR-10 settings; MSE / PSNR / SSIM.
Emphasize held-out CIFAR as the setting that tests a \emph{learned representation}
rather than per-image memorization.

\section{Results}
\subsection{The Decoder Is the Model: Freeze Isolation}
% Strongest, cleanest result first.
Freeze-classical collapses to noise ($\sim$11 dB); freeze-quantum is as good as or
better than the full model. Conclude: the trainable classical decoder is necessary
and sufficient; trained quantum parameters are not.
\subsection{Every Quantum Ablation Ties or Wins (Held-Out CIFAR-10)}
% The core negative result on the trustworthy setting.
Multi-seed held-out table: entangling model ties/trails local-linear classical;
no-entanglement, zz-only, shuffled-pair all indistinguishable. This is the
headline negative result --- report it with error bars.
\subsection{No Generalization from Single Images}
Zero-shot transfer to a new image ($\sim$7.8 dB, below noise floor) shows
single-image runs measure per-image optimization, not representation learning ---
justifying the held-out CIFAR focus.
\subsection{Scaling and Non-Monotonicity}
Bond-dim / qubit / step sweeps are non-monotonic and within noise; HVK is
competitive at 32$\times$32 but degrades to $\sim$19 dB at 256$\times$256 ---
the small observable bottleneck does not scale.

\section{A Circular Benchmark, and How to Catch It}
\label{sec:leakage}
% PROMOTED from supplement to main text — this is genuinely novel and useful.
Present the ``nonlocal advantage'' diagnostic that reported $R^2\!\approx\!1$,
PSNR~$=120$ dB. Show that its regression target is constructed from the exact
pair-product features handed only to the entangling model, so any model given
those columns achieves numerically perfect fit --- entanglement is incidental.
Generalize to a reusable rule: verify no model receives, as a precomputed feature,
a quantity from which the target is linearly recoverable. Frame as a service to
the field, using our own retracted result as the worked example.

\section{Discussion}
\subsection{Why the Physics-Informed Regularizer Was Inert}
Interpret the null Heisenberg-term effect; connect to when physics priors help vs.
when the reconstruction loss already dominates.
\subsection{Dequantization Context}
Relate the tie-with-classical outcome to dequantization theory and the low-entanglement
data regime; explain why entanglement bought nothing here.
\subsection{Design and Evaluation Guidelines}
Concrete bullet-point guidance: resource-match, freeze-isolate, audit leakage,
test on held-out data, report seeds. Position these as the paper's take-home value.
\subsection{Limitations and Scope}
Small circuits, single dataset family, reconstruction (not classification); the
negative result is scoped to this regime and does not preclude advantage elsewhere.

\section{Conclusion}
Restate: a rigorous protocol, a clean negative result, and a transferable leakage
diagnostic advance evaluation methodology in hybrid QML even without a positive
advantage.

\section*{Reproducibility Statement}
% Cheap credibility for a negative-results paper — reviewers weight this heavily.
Point to released code, seeds, resource-matching config, and the exact scripts for
every table. Note the resolved Exp-1 shuffle discrepancy and the retracted circular
benchmark as part of an auditable record.

\section*{Acknowledgements}
Funding, collaborators, student contributions.

\appendix
\section{Exact Circuit Definitions}\label{app:circuits}
Gate sequences for HVK1D / HVK2D / symmetric / IBM-probe; observable lists; the
explicit Heisenberg energy expression from measured observables.
\section{Implementation Details}
Optimizers, learning rates, initializations, training schedules, matched budgets.
\section{Full Ablation Tables}
Complete single-image and multi-image tables with all seeds and error bars.
\section{Extended Leakage Analysis}
Full code excerpt and derivation showing target linear-recoverability from the
leaked features; the corrected task design.
\section{Additional Hardware Results}
IBM depth / CX accounting; shot-noise sweep; transpiler settings.

\end{document}
